\documentclass[11pt]{article}
\usepackage{amsmath,amssymb,amsthm}
\usepackage[margin=1in]{geometry}

\newtheorem{theorem}{Theorem}
\newtheorem{lemma}{Lemma}
\newtheorem{corollary}{Corollary}
\newtheorem{definition}{Definition}

\title{V98: Switching-Balanced Unate Kernels and a Loose-X Label Obstruction}
\author{NC0\_k-Avoid Laboratory}
\date{}

\begin{document}
\maketitle

\paragraph{Setup.}
After essential-support normalization, let a connected local component be
$C:\{0,1\}^n\to\{0,1\}^m$ with $m>n$ and locality at most three.
For every essential unate incidence $(e,v)$ define
$d_{e,v}=0$ when output $e$ is nondecreasing in input $v$, and $d_{e,v}=1$
when it is nonincreasing.

\begin{definition}
The labeled incidence graph is \emph{switching-balanced} when there are
$r_v,q_e\in\{0,1\}$ satisfying
\[
d_{e,v}=r_v\oplus q_e
\]
on every incidence.
\end{definition}

\begin{lemma}
Switching-balance is decidable in linear time in the incidence size.
\end{lemma}

\begin{proof}
The equations are parity constraints on the bipartite input--output incidence
graph. Assign one potential in each connected component and propagate by BFS.
A conflict is exactly a parity-inconsistent incidence cycle. Locality three
makes the incidence size linear in the circuit description.
\end{proof}

\begin{theorem}[Switching theorem]
If the component is switching-balanced, then the coordinate changes
\[
z_v=x_v\oplus r_v,\qquad C'_e(z)=C_e(x)\oplus q_e
\]
make every local gate of $C'$ monotone.
Moreover,
\[
y\in \operatorname{Range}(C)
\quad\Longleftrightarrow\quad
y\oplus q\in\operatorname{Range}(C').
\]
\end{theorem}

\begin{proof}
Input complementation toggles the unate direction by $r_v$ and output
complementation toggles it by $q_e$. Hence the transformed direction is
$d_{e,v}\oplus r_v\oplus q_e=0$. Both coordinate maps are bijections, giving
the range equivalence.
\end{proof}

\begin{corollary}
Every positive-surplus switching-balanced unate $NC^0_3$ component admits
deterministic polynomial-time range avoidance.
\end{corollary}

\begin{proof}
Apply the switching theorem and then the deterministic polynomial-time
monotone-$NC^0_3$ avoider of Kuntewar--Sarma (APPROX/RANDOM 2025) for $m>n$.
Map the missing word back by XOR with $q$.
\end{proof}

\begin{theorem}[Strict V97-kernel family]
For every $N\ge 5$ there is a connected exact-stretch, essential-ternary,
nonmonotone circuit with V97 peeling parameter $\lambda=N$ that is
switching-balanced.
\end{theorem}

\begin{proof}
Use cyclic triple supports $\{i,i+1,i+2\}$, one extra copy of the first support,
and gates
\[
g_e(x)=\operatorname{MAJ}_3(x_v\oplus r_v:v\in S_e),
\qquad r_0=1,\ r_i=0\ (i>0).
\]
Every input has support degree at least three, so no V97 leaf/unary reduction
applies. Gates incident to input zero are decreasing there, so the raw circuit
is nonmonotone. Taking $q_e=0$ satisfies the switching equations.
\end{proof}

\begin{theorem}[Parity loose-X obstruction]
For every $\ell\ge3$ there is a parity-labeled simple loose $X_{2\ell}$ whose
induced edge map is surjective.
\end{theorem}

\begin{proof}
Let the even Berge cycle have vertices $v_0,\ldots,v_{2\ell-1}$. Give edges
$e_0$ and $e_3$ an extra common vertex $w$, and every other edge $e_i$ a private
third vertex $p_i$. For an arbitrary target $y$, set
\[
w=0,\quad v_0=y_0,\quad v_3=y_3,\quad v_i=0\text{ otherwise},
\]
and for $i\notin\{0,3\}$ set
\[
p_i=y_i\oplus v_i\oplus v_{i+1}.
\]
Every parity edge then evaluates to its prescribed target bit.
\end{proof}

\paragraph{Boundary.}
The parity obstruction can be embedded into an exact-stretch connected host
with minimum input degree at least two, so the support-only failure persists
inside the V97 irreducible regime. The host is affine and globally easy; this
is only a falsifier for arbitrary-label loose-X certification. No unrestricted
$NC^0_3$ algorithm, lower-bound transfer, novelty, or P-versus-NP claim is made.

\end{document}
